\section{Exercises}\label{sec:03-propositions-and-proofs:08-exercises:1}

\textbf{Key points.} \texttt{Prop} is the type of statements, and a proof is just an ordinary term of that type (Curry--Howard). \texttt{∧}/\texttt{∨}/\texttt{→}/\texttt{∀}/\texttt{∃} each have their own introduction and elimination rule, mirrored directly by Lean syntax: the anonymous constructor, \texttt{Or.inl}/\texttt{Or.inr}/\texttt{match}, \texttt{fun}/ application, and \texttt{fun}/anonymous-constructor again. \texttt{rfl} proves equality by checking both sides reduce to the same normal form.

\textbf{Socratic questions.}

\begin{enumerate}
\def\labelenumi{\arabic{enumi}.}
\item
  \emph{\texttt{P ∧ Q} and \texttt{P ∨ Q} are both built from two propositions. Why does \texttt{∧} need only one constructor (the anonymous pair) while \texttt{∨} needs two (\texttt{Or.inl}, \texttt{Or.inr})?} A proof of \texttt{P ∧ Q} must supply \emph{both} a proof of \texttt{P} and a proof of \texttt{Q} at once, so one shape suffices. A proof of \texttt{P ∨ Q} supplies \emph{only one} of the two, so the proof itself must record which side was chosen, hence two distinct constructors rather than one.
\item
  \emph{\texttt{rfl} closes \texttt{2 + 2 = 4} immediately. Why can it not, by the same reasoning, close \texttt{∀ n, 0 + n = n}?} \texttt{rfl} checks that both sides already reduce to the same term via unfolding, even with a free variable \texttt{n} in play. \texttt{Nat.add} recurses on its \emph{second} argument, so \texttt{n + 0 = n} actually \textbf{does} close by \texttt{rfl} (it is the base case of the recursion for \texttt{n} itself). But \texttt{0 + n = n} does not. With \texttt{n} unknown, \texttt{0 + n} is stuck on the shape of the second argument and cannot unfold any further without first knowing whether \texttt{n} is \texttt{0} or a successor. That gap between ``true'' and ``reduces to the same term'' (see the discussion of definitional vs. propositional equality in Chapter 5) is exactly why proofs, not mere computation, exist. \texttt{0 + n = n} needs induction (Chapter 4), not \texttt{rfl}.
\item
  \emph{\texttt{∃ p, p > 3 ∧ isPrime p} was proved with a witness and a decided fact about it. What would go wrong trying to prove \texttt{∀ n, n > 0} the same way, by picking one \texttt{n} and checking it?} A single witness proves only that \emph{some} \texttt{n} works. Universal quantification demands the property hold for \emph{every} \texttt{n} at once, which no finite number of individual checks can establish. It needs a genuine argument (an inductive proof, starting in Chapter 4) that works uniformly.
\item
  Prove \texttt{theorem and\_\allowbreak{}comm\_\allowbreak{}ex \{P Q : Prop\} (h : P ∧ Q) : Q ∧ P}.
\item
  Prove \texttt{theorem or\_\allowbreak{}comm\_\allowbreak{}ex \{P Q : Prop\} (h : P ∨ Q) : Q ∨ P} (hint: use \texttt{Or.elim} or pattern matching with \texttt{match}).
\item
  Prove \texttt{theorem exists\_\allowbreak{}gt\_\allowbreak{}zero : ∃ n : Nat, n > 0}.
\end{enumerate}

Solutions: \hyperref[sec:14-appendix-solutions:02-chapter-3:1]{Appendix, Chapter 3}.

